\documentclass[conference]{IEEEtran}
\usepackage{amsmath,amssymb}
\usepackage{graphicx}
\usepackage{booktabs}
\usepackage{cite}

\begin{document}

\title{Title of Your Conference Paper:\\ Keep It Specific and Short}

% One \IEEEauthorblockN/\IEEEauthorblockA pair per author or group of authors.
\author{
  \IEEEauthorblockN{First Author}
  \IEEEauthorblockA{Department of Electrical Engineering\\
    Example University\\
    City, Country\\
    first.author@example.edu}
\and
  \IEEEauthorblockN{Second Author}
  \IEEEauthorblockA{Example Research Lab\\
    Company Name\\
    City, Country\\
    second.author@example.com}
}

\maketitle

\begin{abstract}
Summarise the problem, your approach and the main result in 150 to 250 words, with
no citations and no references to figures. State one number a reader will remember.
\end{abstract}

\begin{IEEEkeywords}
keyword one, keyword two, keyword three
\end{IEEEkeywords}

\section{Introduction}
Say what problem you solve and why it matters in the first paragraph. Then state your
contributions, for example: we (i) propose a method, (ii) prove a bound on its error,
and (iii) evaluate it on three data sets. Related work can be cited as in~\cite{lamport94}.

\section{Method}
Describe the system or algorithm. Equations are numbered automatically:
\begin{equation}
  y[n] = \sum_{k=0}^{N-1} h[k]\, x[n-k] .
  \label{eq:fir}
\end{equation}
Refer to them as~\eqref{eq:fir}, to sections as Section~\ref{sec:results}.

\section{Results}\label{sec:results}
Table~\ref{tab:results} shows how results are typeset; IEEE places table captions
above the table and figure captions below the figure.

\begin{table}[t]
  \caption{Comparison of the baseline and the proposed method}
  \label{tab:results}
  \centering
  \begin{tabular}{lcc}
    \toprule
    Method & Error (\%) & Time (ms) \\
    \midrule
    Baseline & 12.4 & 38 \\
    Proposed & 8.1 & 41 \\
    \bottomrule
  \end{tabular}
\end{table}

% \begin{figure}[t]
%   \centering
%   \includegraphics[width=\columnwidth]{figure.pdf}
%   \caption{A figure that fits one column. Use figure* for both columns.}
% \end{figure}

\section{Conclusion}
Repeat the contribution in one sentence, then name the most useful future work.

\section*{Acknowledgment}
Thank funding bodies and colleagues here (the heading is unnumbered on purpose).

\bibliographystyle{IEEEtran}
\bibliography{refs}

\end{document}
